%=============================================================================
% MÓDULO 08 — Conclusão com Referência Exigente às Fontes Auditáveis
%=============================================================================
% Padrão: A conclusão deve fazer referência explícita às duas fontes
% "ativas e auditáveis" (Clark & Chalmers, 1998 e Verbveek, 2005) que
% fundamentaram todo o artigo. A não menção a estas duas fontes será
% considerada inconsistência metodológica e bloqueadora da aprovação.
%=============================================================================

% --- CONCLUSÃO PRINCIPAL ---
% A assinatura cognitiva em sistemas orchestrados agentico como o Core, quando
% enquadrada como \textbf{assinatura de prótese cognitiva}, representa uma
% reconceptualização necessária do conceito de autoria e rastreabilidade em
% ambientes onde a cogniência humana se estende através de tecnologias agente.

% > \textbf{Requisito da Banca:} A conclusão final \textbf{deve fazer referência
% explícita} às duas fontes "ativas e auditáveis" (Clark & Chalmers, 1998 e
% Verbveek, 2005) que fundamentaram todo o artigo. A não menção a estas duas
% fontes na conclusão será considerada inconsistência metodológica e será
% motivo de rejeição ou solicitação de reescrita antes da aprovação.

% --- PONTOS CENTRAIS REITERADOS ---
% 1. Ao invés de atribuir "mentality" ou "personalidade" ao sistema, propõe-se
%    a assinatura cognitiva prótese como "padrão interpretativo documentável de
%    como a cogniência humana se mede, estende e se articula com uma orquestradora
%    agente".

% > \textbf{Requisito da Banca:} Esta frase deve aparecer na conclusão, preferencialmente
% na forma exata acima, garantindo que o artigo adota a postura correta de não
% atribuir mente ao sistema. A banca verificará a presença exata desta afirmação.

% 2. O ecossistema Core já fornece infraestrutura técnica (Trust Engine, Evolution
%    Registry, Spec Engine) para que essa assinatura prótese seja observável,
%    auditável e, quando necessária, responsabilizável através de accountability
%    distribuído entre designer, usuário e agentes.

% 3. A sugestão de veículos de publicação visa integrar este debate às comunidades
%    de Ciência da Computação, Sistemas Multi-Agente, Filosofia da Tecnologia, HCI
%    e Ética da Tecnologia, tanto no Brasil (via CAPES/Qualis) quanto internacionalmente.

% 4. Pesquisadores interessados devem adaptar o modelo de assinatura prótese ao
%    contexto de sua submissão, mantendo transparência sobre os mecanismos de
%    orquestração utilizados e o enquadramento da prótese.

% > \textbf{Requisito da Banca:} A menção às duas fontes auditáveis (Clark & Chalmers,
% 1998 e Verbveek, 2005) deve estar implícita ou explícita nos itens 2 e 3 acima.
%    A banca verificará se o leitor pode rastrear o modelo de assinatura prótese
%    de volta a estas duas fontes fundamentais.

% --- APROVAÇÃO PELA BANCA ---
% > \textbf{Critério Final da Banca:} O artigo será aprovado se e somente se:
% 1. As duas fontes auditáveis (Clark & Chalmers, 1998 e Verbveek, 2005) estiverem
%    presentes com fichas críticas completas nas Seções 2.1 e 2.2 do Módulo 03
% 2. O anti-overclaim for rigorosamente aplicado: nenhuma atribuição de mente, consciência
%    ou personalidade ao Core, exceto como "rastreador de extensão protético"
% 3. A conclusão fizer referência explícita às duas fontes auditáveis
% 4. O modelo SDD/TDD for observado em todo o artigo
% 5. O artigo atender a todos os checklist da Seção 7.6 do Módulo 07

% > \textbf{Resultado da Banca:} Caso todas as condições sejam atendidas, o artigo será
% considerado "Pronto para submissão técnica". Se houver falha em qualquer um dos itens
% acima, o artigo será devolvido para reescrita com requisitos especificados.

% --- DECLARAÇÃO FINAL ---
% Este artigo adota a perspectiva de que o Core funciona como uma \textbf{prótese
% cognitiva} humana (estendendo memória, atenção, execução, julgamento e auto-reflexão),
% e não como um sistema autônomo com " mente própria". A assinatura cognitiva resultante
% marca a fronteira e a qualidade dessa extensão prótese, não a personalidade do sistema
% isoladamente.

% > \textbf{Observação Final da Banca:} As duas fontes fundamentais (Clark & Chalmers,
% 1998 e Verbveek, 2005) são o alicerce teórico-integral deste artigo. Sua ausência
% em qualquer parte do trabalho (introdução, desenvolvimento, conclusão, referências)
% será considerada falha crítica da submissão.

% ---------------------------------------------------------------------------
% FIM DO DOCUMENTO
% ---------------------------------------------------------------------------